% Section II — Linearized kinetics in entropy coordinates
\section{Linearized kinetics in entropy coordinates}\label{sec:kinetics}

\subsection{Modes, entropy coordinates and the collision operator}\label{sec:modes}

We work with a finite, discretized description of the phonon gas; continuum limits are taken explicitly
in Sec.~\ref{sec:ir}.
\begin{enumerate}
\item[(K1)] A crystal at uniform temperature $T>0$ is described on a finite $\Gamma$-centered mesh of
$N_{\mathbf k}$ wave vectors, invariant under $\mathbf k\mapsto-\mathbf k$; the crystal volume is
$\mathcal V=N_{\mathbf k}V_{\rm cell}$. The modes $\nu=(\mathbf k,s)$, $\nu=1,\dots,n$, have energies
$\varepsilon_\nu=\hbar\omega_\nu>0$ (zero-frequency modes at $\Gamma$ are excluded), group velocities
$\mathbf v_\nu$ and Bose occupations $f^0_\nu=[\exp(\varepsilon_\nu/k_BT)-1]^{-1}$.
\item[(K2)] The occupations obey the linearized Boltzmann equation
$\partial_t\delta f_\nu+\mathbf v_\nu\cdot\nabla\delta f_\nu=-(\Omega\,\delta f)_\nu$, and the collision
operator satisfies detailed balance and the $H$ theorem: with $\mathsf W=\mathrm{diag}(w_\nu)$ and
$w_\nu=[f^0_\nu(1+f^0_\nu)]^{1/2}$, the matrix $\Omega\mathsf W^2$ is symmetric and positive semidefinite.
\end{enumerate}
In the entropy coordinates
\begin{equation}
y_\nu=\frac{\delta f_\nu}{\mathcal V^{1/2}w_\nu}
\label{eq:entropycoords}
\end{equation}
the kinetic equation reads
\begin{equation}
\partial_t\yy+\mathsf V\cdot\nabla\yy=-\Cop\,\yy,\qquad \Cop=\mathsf W^{-1}\Omega\,\mathsf W=\Cop\tr\succeq0,
\label{eq:kinetic}
\end{equation}
with $\mathsf V=\mathrm{diag}(\mathbf v_\nu)$. Full linearized collision operators of this kind are
computed from first principles~\cite{Chaput2013,Fugallo2013,CepellottiMarzari2016}. The name of the
coordinates refers to the entropy density: expanding the
Bose entropy $s_B(f)=k_B[(1+f)\ln(1+f)-f\ln f]$ of each mode to second order gives
\begin{equation}
\delta s=\frac{\delta E}{T}-\frac{k_B}{2}\,|\yy|^2+O(|\yy|^3),\qquad \delta E=\ee\tr\yy,
\label{eq:entropy}
\end{equation}
where $e_\nu=\mathcal V^{-1/2}w_\nu\varepsilon_\nu$, because $s_B'(f^0_\nu)=\varepsilon_\nu/T$ and
$s_B''(f^0_\nu)=-k_B/w_\nu^2$. At fixed energy the Euclidean norm of $\yy$ measures the entropy deficit,
and collisions produce entropy at the rate $k_B\,\yy\tr\Cop\yy\ge0$. Energy conservation in collisions
gives $\Cop\ee=0$.
\begin{enumerate}
\item[(K3)] The energy is the only collision invariant: $\ker\Cop=\mathrm{span}\{\ee\}$.
\end{enumerate}
Hypothesis (K3) excludes exact conservation of crystal momentum (umklapp and other resistive processes
are present). If further invariants exist, everything below holds with $\Hsp$ replaced by the
orthogonal complement of all invariants. We write $\Hsp=\ee^\perp$, $n_{\Hsp}=\dim\Hsp=n-1$ and
$\Pi_{\Hsp}$ for the orthogonal projector onto $\Hsp$; $\Cop$ maps $\Hsp$ into itself and is positive
definite there.

The heat-flux density along the Cartesian direction $i$ is $q_i=\jj^{(i)\mathsf T}\yy$ (a flux per unit
area, in $\mathrm{W\,m^{-2}}$), with the current vector
\begin{equation}
j^{(i)}_\nu=\mathcal V^{-1/2}\,w_\nu\,\varepsilon_\nu\,v_{\nu,i}.
\label{eq:current}
\end{equation}
Time reversal $\nu=(\mathbf k,s)\mapsto\bar\nu=(-\mathbf k,s)$ leaves $\ee$ invariant and reverses
$\jj^{(i)}$, so $\jj^{(i)}\in\Hsp$. The harmonic quantities
\begin{equation}
\rho c=\frac{|\ee|^2}{k_BT^2},\qquad
\Lambda_i=\frac{|\jj^{(i)}|^2}{k_BT^2}=\frac1{\mathcal V}\sum_\nu c_\nu v_{\nu,i}^2,
\label{eq:harmonic}
\end{equation}
are the volumetric heat capacity and a ballistic weight; $c_\nu=\varepsilon_\nu^2w_\nu^2/(k_BT^2)$ is the
heat capacity of mode $\nu$. The first identity follows from the response $\yy=\ee\,\delta T/(k_BT^2)$
of the equilibrium to a temperature change $\delta T$.

\subsection{Dynamic conductivity and memory kernel}\label{sec:response}

Consider the response to a spatially uniform temperature gradient $\nabla T(t)$ switched on at $t=0$
(the zero-wave-vector limit). Write the occupations as the local equilibrium $f^0_\nu(T(\mathbf r,t))$,
with the local temperature fixed by the energy density, plus a deviation whose entropy coordinates
$\yy$ lie in $\Hsp$. The streaming of the local equilibrium,
$\mathbf v_\nu\cdot\nabla f^0_\nu=w_\nu^2\varepsilon_\nu\,\mathbf v_\nu\cdot\nabla T/(k_BT^2)$, acts as a
source with entropy coordinates $\sum_k\jj^{(k)}\partial_kT/(k_BT^2)\in\Hsp$; the time derivative of the
local equilibrium is proportional to $\ee$ and drops out on projection onto $\Hsp$ (its projection onto
$\ee$ is the energy balance). For a spatially uniform deviation,
\begin{equation}
\partial_t\yy=-\Cop\yy-\frac{1}{k_BT^2}\sum_k\jj^{(k)}\,\partial_kT(t),\qquad \yy(0)=0.
\label{eq:drive}
\end{equation}
Its solution gives the heat flux in the form of a constitutive law with memory~\cite{GurtinPipkin1968},
\begin{equation}
q_i(t)=-\int_0^t\kappa_{ik}(t-t')\,\partial_kT(t')\,dt',\qquad
\kappa_{ik}(t)=\frac{\Phi_{ik}(t)}{k_BT^2},
\label{eq:memorylaw}
\end{equation}
with the current autocorrelation function $\Phi_{ik}(t)=\jj^{(i)\mathsf T}e^{-\Cop t}\jj^{(k)}$. Under the
time convention $e^{+i\omega t}$ of this paper, the Laplace variable is $p$ and the harmonic regime is
$p=i\omega$. The transform of Eq.~\eqref{eq:memorylaw} defines the dynamic conductivity
\begin{equation}
\lambda_{ik}(p)=\frac{K_{ik}(p)}{k_BT^2},\qquad K_{ik}(p)=\jj^{(i)\mathsf T}(p+\Cop)^{-1}\jj^{(k)}.
\label{eq:dynamicK}
\end{equation}
Because the currents lie in $\Hsp$, where $\Cop$ is invertible, $K_{ik}$ is analytic at $p=0$ and the
dc conductivity is $\lambda_{ik}=K_{ik}(0)/(k_BT^2)$ with $K_{ik}(0)=\jj^{(i)\mathsf T}\Cop^+\jj^{(k)}$;
$\Cop^+$ is the Moore--Penrose pseudoinverse. A current with a component along an invariant would add a
term proportional to $1/p$, and the pseudoinverse must not be used to hide it. From now on we fix one
direction and drop the index $i$: $\jj$, $K(p)$, $\Phi(t)$, $\lambda(p)$ and $\lambda=\lambda(0)$.

With the eigenpairs $(\Gamma_k,\mathbf u_k)$ of $\Cop$ restricted to $\Hsp$ ($\Gamma_k>0$; the relaxons
of Ref.~\cite{CepellottiMarzari2016}),
\begin{equation}
K(p)=\sum_{k=1}^{n-1}\frac{(\mathbf u_k\tr\jj)^2}{p+\Gamma_k}.
\label{eq:stieltjes}
\end{equation}
$K$ is the Stieltjes transform of the spectral measure of $\Cop|_{\Hsp}$ in the state $\jj$; its total
mass is $|\jj|^2=\Phi(0)$.

\subsection{The memory moment}\label{sec:memory}

\begin{definition}[Moments]\label{def:moments}
Under (K1)--(K3) and for $\jj\in\Hsp\setminus\{0\}$, let $\xx=\Cop^+\jj$ (the unique solution in $\Hsp$
of $\Cop\xx=\jj$), and define
\begin{equation}
\begin{gathered}
K_0=K(0)=\jj\tr\xx,\qquad N=|\xx|^2,\\
\tmem=-\frac{K'(0)}{K(0)},\qquad \tstat=\frac{K_0}{|\jj|^2}.
\end{gathered}
\label{eq:moments}
\end{equation}
\end{definition}

\begin{lemma}[Equivalent forms of the memory moment; \status{proved}]\label{lem:forms}
Under the hypotheses of Definition~\ref{def:moments},
\begin{equation}
\tmem=\frac{N}{K_0}=\frac{|\Cop^+\jj|^2}{\jj\tr\Cop^+\jj}
=\frac{\int_0^\infty t\,\Phi(t)\,dt}{\int_0^\infty\Phi(t)\,dt}
=\frac{\sum_k\lambda_k\Gamma_k^{-1}}{\sum_k\lambda_k},
\label{eq:forms}
\end{equation}
where $\lambda_k=(\mathbf u_k\tr\jj)^2/(k_BT^2\,\Gamma_k)$ is the contribution of relaxon $k$ to
$\lambda$. In particular $\Gamma_{\max}^{-1}\le\tmem\le\Gamma_{\min}^{-1}$, where $\Gamma_{\min}$ and
$\Gamma_{\max}$ are the extreme eigenvalues of $\Cop|_{\Hsp}$.
\end{lemma}

\begin{proof}
On $\Hsp$, $(p+\Cop)^{-1}$ is analytic at $p=0$ with derivative $-(p+\Cop)^{-2}$; hence
$-K'(0)=\jj\tr(\Cop^+)^2\jj=|\Cop^+\jj|^2$ because $\Cop^+$ is symmetric. In time,
$\int_0^\infty e^{-\Cop t}dt=\Cop^{+}$ and $\int_0^\infty t\,e^{-\Cop t}dt=(\Cop^{+})^2$ on $\Hsp$. The
last form follows from Eq.~\eqref{eq:stieltjes}: $K_0=\sum_k(\mathbf u_k\tr\jj)^2\Gamma_k^{-1}$ and
$N=\sum_k(\mathbf u_k\tr\jj)^2\Gamma_k^{-2}$. A weighted mean of the $\Gamma_k^{-1}$ lies between their
extremes.
\end{proof}

Thus $\tmem$ is the first moment of the normalized memory kernel $\kappa(t)/\lambda$ and the
conductivity-weighted mean relaxon lifetime. In the RTA, where $\Cop$ is replaced by
$\mathrm{diag}(r_\nu)$ (an operator that does not conserve energy),
$\tmem^{\rm RTA}=\sum_\nu j_\nu^2r_\nu^{-2}/\sum_\nu j_\nu^2r_\nu^{-1}$ and
$\tstat^{\rm RTA}=\sum_\nu j_\nu^2r_\nu^{-1}/\sum_\nu j_\nu^2$.

\subsection{The Cattaneo time of a conserving closure}\label{sec:cattaneo}

For spatially uniform fields, Cattaneo's law $\tau_R\partial_tq+q=-\lambda\,\partial_iT$ is
Eq.~\eqref{eq:memorylaw} with the exponential kernel $\kappa_C(t)=(\lambda/\tau_R)e^{-t/\tau_R}$, that
is, with the single pole $\lambda_C(p)=\lambda/(1+p\tau_R)$~\cite{JosephPreziosi1989,GurtinPipkin1968}.

\begin{definition}[Cattaneo time]\label{def:cattaneo}
The Cattaneo time of a conserving operator $\Cop$ (for a given direction) is the relaxation time of the
single pole that matches $\lambda(p)$ and $\lambda'(p)$ at $p=0$. Since $\lambda_C(0)=\lambda$ and
$\lambda_C'(0)=-\lambda\tau_R$, it equals $\tmem$.
\end{definition}

The match is exact for all $p$ if and only if $\Phi$ is a single exponential, that is, if and only if
$\jj$ is an eigenvector of $\Cop$. The matched pole reproduces the slope of $\lambda(p)$ at $p=0$, not
the response in a finite frequency band; Theorem~\ref{thm:psd}(v) below shows that the two can be
unrelated.

The grey conserving closure illustrates Definition~\ref{def:cattaneo}. In the Callaway
model~\cite{Callaway1959} with a common linear isotropic dispersion $\varepsilon=\hbar v|\mathbf k|$
and constant times $\tau_N,\tau_R>0$, the linearized operator in entropy coordinates is
\begin{equation}
\Cop_{\rm Cal}=\tau_N^{-1}(I-\Pi_{\rm N})+\tau_R^{-1}\,\Pi_{\Hsp},
\label{eq:callaway}
\end{equation}
where $\Pi_{\rm N}$ is the orthogonal projector onto the span of $\ee$ and of the three vectors
$\mathbf P_x$, $\mathbf P_y$, $\mathbf P_z$, with $P_{i,\nu}=\mathcal V^{-1/2}w_\nu\hbar k_{\nu,i}$,
which represent crystal momentum: the displaced
equilibrium that matches the energy and momentum of a state $\yy$ is its orthogonal projection onto that
span. We take the mesh symmetric under the cubic point group, so that $\ee,\mathbf P_x,\mathbf P_y,\mathbf P_z$ are
mutually orthogonal.

\begin{proposition}[Grey conserving closure; \status{derived under assumptions}]\label{prop:callaway}
Under the assumptions just stated, $\jj^{(i)}=v^2\mathbf P_i$ and $\Cop_{\rm Cal}\jj^{(i)}=\jj^{(i)}/\tau_R$.
Hence $\lambda(p)=\lambda/(1+p\tau_R)$ exactly, with $\lambda=\rho c\,v^2\tau_R/3$, and
$\tmem=\tstat=\tau_R$. The mode relaxation rates are
$(\Cop_{\rm Cal})_{\nu\nu}=\tau_N^{-1}+\tau_R^{-1}-O(n^{-1})$.
\end{proposition}

\begin{proof}
For $\varepsilon=\hbar v|\mathbf k|$, $\varepsilon_\nu v_{\nu,i}=\hbar v^2k_{\nu,i}$, so
$\jj^{(i)}=v^2\mathbf P_i\in\mathrm{range}\,\Pi_{\rm N}\cap\Hsp$, and Eq.~\eqref{eq:callaway} gives
$\Cop_{\rm Cal}\jj^{(i)}=\jj^{(i)}/\tau_R$. Then $K(p)=|\jj|^2/(p+\tau_R^{-1})$, and
$\lambda=|\jj|^2\tau_R/(k_BT^2)=\Lambda_i\tau_R=\rho c\,v^2\tau_R/3$ by Eq.~\eqref{eq:harmonic} and cubic
symmetry. The diagonal elements of the projectors are $e_\nu^2/|\ee|^2+\sum_iP_{i,\nu}^2/|\mathbf P_i|^2$
and $e_\nu^2/|\ee|^2$, both $O(n^{-1})$.
\end{proof}

Eliminating the deviatoric second moment of the same model under $|p|\tau_N\ll1$ and
$v\tau_N/L_{\rm v}\ll1$, with $L_{\rm v}$ the scale of spatial variation, gives the conserving GK
closure~\cite{GuoWang2022,Sendra2021,Sendra2022}
\begin{equation}
\tau_R\partial_t\mathbf q+\mathbf q=-\lambda\nabla T+\ell^2\big[\nabla^2\mathbf q+\beta\,\nabla(\nabla\cdot\mathbf q)\big],
\label{eq:gk}
\end{equation}
with $\beta=1/3$, $\ell^2=v^2\tau_c\tau_R/5$ and $\tau_c=(\tau_N^{-1}+\tau_R^{-1})^{-1}$; the original
GK equation has $\beta=2$ and $\tau_N$ in place of $\tau_c$~\cite{GuyerKrumhansl1966b}. Its relaxation
time is the $\tau_R$ of Proposition~\ref{prop:callaway}, i.e., $\tmem$. The mode lifetimes, however, are
close to $\tau_c$, not to $\tau_R$: lifetimes alone do not give the Cattaneo time even in the grey
model, whereas within this two-parameter family the pair $(\tau_c,\lambda)$ fixes $(\tau_N,\tau_R)$.
Section~\ref{sec:identifiability} asks what remains of this when the operator is not restricted to a
two-parameter family.

\subsection{Static data and operator classes}\label{sec:classes}

\begin{definition}[Static data]\label{def:data}
For $\Cop$ satisfying (K2)--(K3), the static data are $\mathcal D(\Cop)=(\mathbf r,\ker\Cop,K_0)$, where
$\mathbf r=\mathrm{diag}\,\Cop$ are the mode relaxation rates (lifetimes $\tau_\nu=1/r_\nu$). The
harmonic vectors $\ee$ and $\jj$ are fixed. For several directions, $K_0$ is replaced by the dc tensor.
\end{definition}

Under (K3) all $r_\nu>0$: a vanishing diagonal element of a positive-semidefinite matrix forces the
corresponding row to vanish, which would put a unit vector in $\ker\Cop=\mathrm{span}\{\ee\}$, whereas
$\ee$ has no zero component. The rates $r_\nu$ include the self-coupling of repeated slots in an event;
they are not, in general, the self-energy linewidths of RTA codes (see Lemma~\ref{lem:rta}).

\begin{definition}[Operator classes]\label{def:classes}
(P)~The positive-semidefinite class is
$\mathfrak P(\ee)=\{\Cop\in\mathbb R^{n\times n}:\Cop=\Cop\tr\succeq0,\ \ker\Cop=\mathrm{span}\{\ee\}\}$.
(E)~Let $Z=(\bm\zeta_\gamma)_{\gamma\in\mathcal E}$ be a fixed finite family of vectors in $\Hsp$ that
spans $\Hsp$. For rates $g\in[0,\infty)^{\mathcal E}$ put
\begin{equation}
\Cop(g)=\sum_{\gamma\in\mathcal E}g_\gamma\,\bm\zeta_\gamma\bm\zeta_\gamma\tr .
\label{eq:cone}
\end{equation}
The rates are admissible if $\ker\Cop(g)=\mathrm{span}\{\ee\}$, equivalently if
$\{\bm\zeta_\gamma:g_\gamma>0\}$ spans $\Hsp$. The fixed event cone is
$\mathfrak E(Z)=\{\Cop(g):g\ \text{admissible}\}\subset\mathfrak P(\ee)$.
\end{definition}

Three-phonon kinetics has the form~\eqref{eq:cone}~\cite{Fugallo2013}. For an event
$\gamma:\ \mathsf p\to\mathsf a+\mathsf b$ allowed by crystal-momentum conservation (normal or umklapp)
and by energy conservation on the mesh, let $\mathbf s_\gamma=-\bm\delta_{\mathsf p}+\bm\delta_{\mathsf a}+\bm\delta_{\mathsf b}$
be its stoichiometric vector ($-\bm\delta_{\mathsf p}+2\bm\delta_{\mathsf a}$ for a repeated daughter).
Linearizing the net rate $G_\gamma[f_{\mathsf p}(1+f_{\mathsf a})(1+f_{\mathsf b})-(1+f_{\mathsf p})f_{\mathsf a}f_{\mathsf b}]$
about $f^0$ and using detailed balance gives $-G_\gamma P_\gamma\,\mathbf s_\gamma\tr\mathsf W^{-2}\delta\mathbf f$,
with $P_\gamma=f^0_{\mathsf p}(1+f^0_{\mathsf a})(1+f^0_{\mathsf b})$, so that Eq.~\eqref{eq:kinetic} holds
with
\begin{equation}
\bm\zeta_\gamma=\mathsf W^{-1}\mathbf s_\gamma,\qquad g_\gamma=G_\gamma P_\gamma\ge0 .
\label{eq:events}
\end{equation}
Energy conservation $\mathbf s_\gamma\tr\bm\varepsilon=0$ gives $\bm\zeta_\gamma\in\Hsp$. The event set,
hence $Z$, is fixed by the harmonic spectrum; the rates $g_\gamma$ are fixed by the anharmonic
amplitudes, which are the unknowns of the inverse problem considered here. On a finite mesh, exact
resonance is exceptional; with a broadened resonance condition the accepted events must be given an
energy-conserving stoichiometry, which is a declared approximation (Supplemental Material, Sec.~S1).
The class statements below concern any fixed finite $Z$.
For given data $(\mathbf r^\star,K^\star)$ the feasible set is
\begin{equation}
\begin{split}
\mathcal F(\mathbf r^\star,K^\star)=\{g\ \text{admissible}:\ &\mathrm{diag}\,\Cop(g)=\mathbf r^\star,\\
&K_0(g)=K^\star\}.
\end{split}
\label{eq:feasible}
\end{equation}

Compressing the kinetics onto fewer variables is a separate question. Block observables
$\mathbf N=\mathsf B\yy$ (for instance the populations of degenerate blocks) obey a closed equation
$\dot{\mathbf N}=-\mathsf G\mathbf N$ for every initial state if and only if $\mathsf B\Cop=\mathsf G\mathsf B$,
the exact-lumpability condition~\cite{WeiKuo1969,Aoki1968,LiRabitz1989,Buchholz1994}; otherwise the reduced
dynamics contains memory terms of projection type~\cite{Zwanzig1961,Mori1965}. We use this known
criterion only as background.
